\section{Introduction}
\label{sec:introduction}
Feynman integrals arise in higher-order perturbative corrections to
scattering amplitudes and particle decay amplitudes. In QCD, these
corrections involve quark and gluon loops and contribute to predictions
for cross sections and decay rates. Calculating a loop contribution
requires integration over the independent internal momenta that remain
after imposing momentum conservation~\cite{Weinzierl2010}.

The scalar integrals that appear in an amplitude have the general form
\begin{equation}
 \mathcal J=
 \left(\prod_{k=1}^{L}\int
       \frac{\mathrm{d}^{d}\ell_k}{\mathrm{i}\pi^{d/2}}\right)
 \frac{\mathcal N(\boldsymbol\ell,\boldsymbol p)}
      {\prod_{i=1}^{n_{\mathrm{phys}}}D_i^{a_i}},
 \qquad D_i=q_i^2-m_i^2+\mathrm{i}0.
 \label{eq:general-integral}
\end{equation}
Here $L$ is the number of loops and
$\boldsymbol\ell=(\ell_1,\ldots,\ell_L)$ is the list of independent loop
momenta over which we integrate. The number of independent external
momenta is $E$, and we keep these momenta,
$\boldsymbol p=(p_1,\ldots,p_E)$, fixed during integration. The numerator
$\mathcal N$ is a polynomial obtained from the spin and tensor algebra;
we leave its mass dependence implicit in the notation.
The number of physical denominators is $n_{\mathrm{phys}}$, and each
$D_i$ corresponds to an internal particle line with momentum $q_i$ and
mass $m_i$. The positive integer $a_i$ is the power of $D_i$ in the
denominator, and $q_i$ is a fixed linear combination of loop and external
momenta, with $q_i^2=q_i\cdot q_i$ denoting the Lorentz scalar product.

The integration dimension is $d=4-2\epsilon$, where $\epsilon$ is the
regulator used to define divergent integrals, so the result depends on
$\epsilon$ as well as the external momenta and masses. The factor
$1/(\mathrm{i}\pi^{d/2})$ is our normalization convention for each
integration measure, with $\mathrm{i}^2=-1$. In each denominator,
$\mathrm{i}0$ is an infinitesimal positive imaginary term that specifies
how to integrate past a propagator pole, where $q_i^2=m_i^2$.

Higher loop orders produce collections of integrals with related numerators
and denominators and algebraic identities can be used to express them in
terms of an independent set of master integrals. Kira automates this
reduction~\cite{KiraOriginal}, after which the calculation requires the
values of the master integrals.
To explain how the reduction works and where machine learning (ML) could
speed it up, we first define the notation used in this paper, then derive
the identities and describe how Kira solves the resulting equations.

\section{The mathematical form of Feynman integrals}
\label{sec:integral-form}
To label the integrals related to Eq.~\eqref{eq:general-integral}, we
specify their momentum dependence and write the numerator and denominator
powers as a vector of integers. This vector identifies each integral as
an unknown in the reduction equations.

\subsection{Loop momenta, propagators and regularization}
After choosing how the loop and external momenta flow through the internal
lines, we write each physical denominator as
\begin{equation}
 D_i=q_i^2-m_i^2+\mathrm{i}0,
 \qquad
 q_i=\sum_{k=1}^{L}\alpha_{ik}\ell_k
       +\sum_{\beta=1}^{E}\eta_{i\beta}p_\beta .
 \label{eq:propagator}
\end{equation}
The constants $\alpha_{ik}$ and $\eta_{i\beta}$ specify which loop and
external momenta make up $q_i$, with indices $k,j$ for loop momenta and
$\beta,\gamma$ for external momenta. Squaring $q_i$ gives the scalar
products $\ell_k\cdot\ell_j$, $\ell_k\cdot p_\beta$ and
$p_\beta\cdot p_\gamma$, and those containing loop momenta vary during
integration. The vector $\boldsymbol x$ contains the independent external
scalar products and squared masses, which are the kinematic parameters
we keep fixed during integration.

Loop integrals can diverge at large momenta or when massless momenta become
small or parallel. Dimensional regularization defines these integrals at
values of $d$ for which they converge and extends the result by analytic
continuation to $d=4-2\epsilon$, using additional regulators where
necessary. In an expansion around $\epsilon=0$, the divergences appear
as negative powers of $\epsilon$, which
renormalization and infrared cancellation in the observable address later
in the calculation~\cite{Weinzierl2010}. Having specified the regulated
measure and momentum dependence, we can now use the powers of the
numerator and denominator factors to label individual integrals.

\subsection{A common notation for an integral family}
We add auxiliary expressions to the physical denominators so that we can
write every required loop-momentum scalar product in terms of the list
$D_1,\ldots,D_N$ and the fixed parameters $\boldsymbol x$. Here $N$ is
the number of expressions in the list, and the auxiliary expressions
allow us to represent numerator factors. Writing the numerator in terms
of this list expresses $\mathcal J$ as a linear combination of integrals
of the form
\begin{equation}
 I_{\boldsymbol a}(\boldsymbol x,\epsilon)
 =\int_{\ell}\prod_{i=1}^{N}D_i^{-a_i},
 \qquad
 \int_{\ell}\equiv
 \prod_{k=1}^{L}\int\frac{\mathrm{d}^{d}\ell_k}{\mathrm{i}\pi^{d/2}},
 \qquad \boldsymbol a=(a_1,\ldots,a_N)\in\mathbb Z^N.
 \label{eq:family}
\end{equation}
Here $\int_\ell$ is shorthand for the loop integration measure and
$\boldsymbol a$ is the vector of $N$ integer powers, with $\mathbb Z$
denoting the set of integers. These powers can be positive, zero or
negative, so we can
write numerator and denominator factors in the equivalent form
\begin{equation}
 I_{\boldsymbol a}(\boldsymbol x,\epsilon)
 =\int_\ell
 \frac{\prod_{i:a_i<0}D_i^{-a_i}}
      {\prod_{i:a_i>0}D_i^{a_i}}.
 \label{eq:index-signs}
\end{equation}
A positive index places $D_i^{a_i}$ in the denominator, a negative index
places $D_i^{-a_i}$ in the numerator, and a zero index removes the factor,
with an empty product equal to one. Auxiliary expressions used only as
numerators therefore have nonpositive indices. An irreducible scalar
product is a numerator scalar product that cannot be written as a linear
combination of the physical denominators and terms depending only on
$\boldsymbol x$~\cite{Weinzierl2022}.

An integral family contains the integrals obtained by varying
$\boldsymbol a$ for a fixed list $D_1,\ldots,D_N$ and fixed kinematics.
Within this family, the target list specifies the integrals that the
amplitude requires, and grouping family members by their positive indices
organizes them into sectors.

\subsection{Sectors and the organization of related integrals}
A sector is a set of integrals with the same positive-index positions,
so they contain the same denominators. We assign a label to each sector
using
\begin{equation}
 \theta_i(\boldsymbol a)=
 \begin{cases}1,&a_i>0,\\0,&a_i\leq0,\end{cases}
 \qquad
 S(\boldsymbol a)=\sum_{i=1}^{N}2^{i-1}\theta_i(\boldsymbol a).
 \label{eq:sector}
\end{equation}
The indicator $\theta_i$ is one when denominator $i$ is present and zero
otherwise, and $S$ is an integer that identifies this set of denominators.
Integrals with different powers of the same denominators belong to the
same sector because $S$ depends only on which indices are positive.
Removing one or more denominators gives a subsector containing a proper
subset of the original sector's denominators.

Scaleless sectors contain no nonzero mass or kinematic scale and therefore
vanish in dimensional regularization, while changes of loop integration
variables can yield symmetry relations between integrals. Kira uses these
zero sectors and symmetry mappings to reduce the number of
unknowns~\cite{KiraOriginal}. We can then use the sector and index labels
to write the identities that relate the remaining family members.

\section{From integration to algebraic reduction}
\label{sec:ibp-reduction}
To express the targets in terms of master integrals, we generate
relations between family members through integration by parts and solve
the resulting equations. Evaluating the master integrals then gives
the target integral values needed to calculate the amplitude.

\subsection{How integration by parts produces relations}
For differentiable functions $f(z)$ and $g(z)$ of a real variable $z$, the
product rule gives the integration identity between endpoints $z_0,z_1$:
\begin{equation}
 \int_{z_0}^{z_1} f'(z)g(z)\,\mathrm{d}z
 =[f(z)g(z)]_{z_0}^{z_1}
 -\int_{z_0}^{z_1} f(z)g'(z)\,\mathrm{d}z.
 \label{eq:ordinary-ibp}
\end{equation}
Primes denote derivatives with respect to $z$, and the brackets denote the
endpoint difference $f(z_1)g(z_1)-f(z_0)g(z_0)$. To apply this product rule
to loop-momentum derivatives, we derive the integration-by-parts (IBP)
identity where the regulated integral has a vanishing boundary term and then
extend it by analytic continuation~\cite{ChetyrkinTkachov,Weinzierl2022}.
For a loop momentum $\ell_k$ and a vector
$v\in\{\ell_1,\ldots,\ell_L,p_1,\ldots,p_E\}$, the identity is
\begin{equation}
 0=\int_{\ell}\frac{\partial}{\partial\ell_k^\mu}
       \left(v^\mu\prod_{i=1}^{N}D_i^{-a_i}\right).
 \label{eq:ibp}
\end{equation}
The index $\mu$ is a spacetime index, and a repeated Lorentz index means
that we sum over components using the spacetime metric. The derivative
acts on both $v^\mu$ and the factors $D_i^{-a_i}$, so the product rule gives
\begin{equation}
 0=\int_{\ell}\left[
 \frac{\partial v^\mu}{\partial\ell_k^\mu}
 -\sum_{i=1}^{N}a_i\,
   \frac{v^\mu(\partial D_i/\partial\ell_k^\mu)}{D_i}
 \right]\prod_{j=1}^{N}D_j^{-a_j}.
 \label{eq:ibp-expanded}
\end{equation}
The first term equals $d$ for $v=\ell_k$ and vanishes when $v$ is an
external momentum or another independent loop momentum. Differentiating
an inverse propagator produces scalar products, which we express through
the family basis to convert each term into an integral with shifted indices:
\begin{equation}
 \sum_{j=1}^{n_\lambda}
 c_{\lambda j}(\boldsymbol a;\boldsymbol x,\epsilon)
 I_{\boldsymbol a+\boldsymbol\Delta_{\lambda j}}
 (\boldsymbol x,\epsilon)=0.
 \label{eq:ibp-linear}
\end{equation}
Here $\lambda$ identifies the chosen pair $(k,v)$, $n_\lambda$ is the number of
terms, and $\boldsymbol\Delta_{\lambda j}$ is the integer change in the
index vector for term $j$. Differentiation gives the coefficient
$c_{\lambda j}$, which depends on the integral indices and the parameters
$\boldsymbol x,\epsilon$~\cite{ChetyrkinTkachov}.

The $L$ choices of differentiated momentum and $L+E$ choices of $v$ give
$L(L+E)$ basic identities at each index vector, some of which may depend
linearly on the others. Further relations follow from Lorentz invariance
because a common rotation or Lorentz boost of all external momenta
preserves a scalar integral, giving the infinitesimal identity
\begin{equation}
 \sum_{\beta=1}^{E}\left(
 p_\beta^\nu\frac{\partial}{\partial p_{\beta\mu}}
 -p_\beta^\mu\frac{\partial}{\partial p_{\beta\nu}}
 \right) I_{\boldsymbol a}=0.
 \label{eq:li}
\end{equation}
Here $\mu,\nu$ are free spacetime indices, meaning that the identity holds
for each pair of components, and the derivatives act on the
external-momentum dependence of the integral. Contracting these indices
with external momenta gives scalar relations~\cite{GehrmannRemiddi}
that join the IBP and symmetry identities to form the reduction system.

\subsection{From relations to master integrals}
Let $n_I$ be the number of distinct integrals in the generated equations
and $n_{\mathrm{eq}}$ the number of equations. We arrange these integrals
in a column vector so that the system takes the matrix form
\begin{equation}
 A(\boldsymbol x,\epsilon)\,\boldsymbol I=0,
 \qquad
 \boldsymbol I=(I_{\boldsymbol a^{(1)}},\ldots,
 I_{\boldsymbol a^{(n_I)}})^{\mathsf T}.
 \label{eq:matrix}
\end{equation}
The superscript $\mathsf T$ denotes transpose, so $\boldsymbol I$ is a
column vector of integrals whose dependence on $\boldsymbol x,\epsilon$
we leave implicit here. The matrix $A$ has $n_{\mathrm{eq}}$ rows and
$n_I$ columns, with each entry $A_{ij}$ giving the coefficient of integral
$j$ in equation $i$. Solving this system expresses dependent integrals in
terms of an independent set. This set is a master basis if the reduction
relations express every required integral as a linear combination of its
members. For the standard class of Feynman integrals at fixed graph,
this basis contains a finite number
of masters~\cite{SmirnovPetukhov}. We write them as
$M_1,\ldots,M_{n_M}$, where $n_M$ is the number of master integrals, so
the reduction formula is
\begin{equation}
 I_{\boldsymbol a}(\boldsymbol x,\epsilon)
 =\sum_{m=1}^{n_M}C_{\boldsymbol a m}(\boldsymbol x,\epsilon)
 M_m(\boldsymbol x,\epsilon).
 \label{eq:reduction}
\end{equation}
Here $C_{\boldsymbol a m}$ is the coefficient of master $M_m$ in the
reduction of target $I_{\boldsymbol a}$. The solver obtains it by adding,
multiplying and dividing the identity coefficients, so it is a rational
function, meaning a ratio of polynomials in $\boldsymbol x,\epsilon$. Different
independent sets can define different master bases, and an insufficient
equation set can leave additional integrals unresolved. Establishing the
master basis therefore requires checking that the generated equations
are enough to reduce the required integrals.

Equation~\eqref{eq:reduction} gives the coefficients multiplying the
masters, whose evaluation requires a further calculation, for example by
solving differential equations with known values or limits as boundary
conditions~\cite{GehrmannRemiddi}. Kira provides the reduction rules
needed for this calculation by generating and processing the identities
through the procedure described next.

\section{How Kira performs the reduction}
\label{sec:kira-procedure}
Kira implements the Laporta approach by generating a finite set of
identities and eliminating integrals according to an
ordering~\cite{Laporta}. To obtain the coefficients in
Eq.~\eqref{eq:reduction}, Kira selects seeds, removes redundant
equations and solves the retained system. We examine how Kira performs
these steps to specify the seed and elimination choices that the
proposed models would make.

\subsection{Kira's input and seed generation}
The user supplies the families in Eq.~\eqref{eq:family}, relations among
the parameters $\boldsymbol x$, the target list, and bounds on equation
generation so that Kira can initialize the families and sectors. During
this initialization, Kira detects scaleless sectors and identifies symmetry
relations and mappings between compatible families~\cite{KiraOriginal}.

A seed specifies the index vector $\boldsymbol a$ at which Kira evaluates
Eq.~\eqref{eq:ibp-linear} to generate several derivative identities, and
the seed set may include integrals beyond the target list. The shifts
$\boldsymbol\Delta_{\lambda j}$ in these identities can introduce further
integrals outside the seed set, so the resulting system depends on both
the chosen seeds and the terms generated from them. To set limits on the
indices of the seeds, Kira uses the quantities
\begin{align}
 t(\boldsymbol a)&=\sum_{i=1}^{N}\theta_i(\boldsymbol a),
 &r(\boldsymbol a)&=\sum_{a_i>0}a_i,\\
 s(\boldsymbol a)&=-\sum_{a_i<0}a_i,
 &n_{\mathrm{dot}}(\boldsymbol a)&=r(\boldsymbol a)-t(\boldsymbol a).
 \label{eq:complexities}
\end{align}
Here $t$ is the number of distinct denominators, $r$ is the sum of their
powers, and $s$ is the sum of the numerator powers. The quantity
$n_{\mathrm{dot}}$ is the sum of the extra denominator powers beyond
one for each denominator. We use this symbol to keep $d$ reserved for the
integration dimension, and we restrict equation generation by specifying
the allowed sectors and limits
\begin{equation}
 r\leq r_{\max},\qquad s\leq s_{\max},\qquad
 n_{\mathrm{dot}}\leq n_{\mathrm{dot},\max},
 \label{eq:seed-bounds}
\end{equation}
where $r_{\max}$, $s_{\max}$ and $n_{\mathrm{dot},\max}$ are the upper
limits on the corresponding quantities. The user sets $r_{\max}$ and
$s_{\max}$ and may also set $n_{\mathrm{dot},\max}$. Because these choices determine which seeds
contribute equations, a reduction run must establish whether the bounds
supply enough relations for the targets~\cite{Kira12}.

Kira 3 changes seeding in symmetry-related sectors and provides the optional
\texttt{truncate\_sp} setting to restrict numerator seeding in subsectors,
with parameter choices that require checks against the expected master
basis~\cite{Kira3}. The generated system can still contain redundant
equations, which elimination identifies as it processes the relations
into reduction rules.

\subsection{Integral ordering, sector order and forward elimination}
Elimination solves an identity for one integral and substitutes the result
into other identities, using an integral ordering to select which unknown
to eliminate from each equation. Kira defines these orderings through
sector labels and the powers of the numerator and denominator factors,
with choices that prioritize numerator powers or repeated denominators
and thereby affect the master basis and
runtime~\cite{Kira12}.

To describe this procedure, write $I_j=I_{\boldsymbol a^{(j)}}$ for the
integrals in Eq.~\eqref{eq:matrix} and choose the ordering
$I_1\succ I_2\succ\cdots\succ I_{n_I}$, where $\succ$ gives priority
for elimination. In a row whose first nonzero term contains $I_k$, let
$b_j$ denote the coefficient of $I_j$, so choosing the nonzero coefficient
$b_k$ as the pivot gives
\begin{equation}
 b_k I_k+\sum_{j>k}b_j I_j=0
 \quad\Longrightarrow\quad
 I_k=-\sum_{j>k}\frac{b_j}{b_k}I_j.
 \label{eq:pivot}
\end{equation}
The coefficients $b_j$ depend on $\boldsymbol x,\epsilon$, and substituting
this expression into the other rows removes their dependence on $I_k$.
Repeating this step gives forward elimination, which builds a triangular
system in which each pivot integral depends on integrals later in the
ordering. Back substitution then replaces the remaining pivot integrals
using their own rules, completing the Laporta solution~\cite{Laporta}.

Sector scheduling sets the order in which the solver processes groups of
integrals or equations, taking into account which other integrals their
equations contain. Ordinary IBP identities for a sector involve that
sector and its subsectors because nonpositive indices remain nonpositive
when applying the derivative. Equations for one sector can therefore
depend on integrals from its subsectors. The Kira 3 description starts
equation generation with sectors earlier in its ordering and inserts
previously obtained pivot relations into new equations~\cite{Kira3},
with the order of processing depending on the stage of the reduction.

When several rows share the same first integral, the choice of pivot row
changes the substitutions that follow. Choosing row $\rho$ to eliminate
column $k$, for example, updates row $i$ by
\begin{equation}
 A_{ij}\ \longleftarrow\ A_{ij}
 -\frac{A_{ik}}{A_{\rho k}}A_{\rho j},
 \qquad A_{\rho k}\ne0.
 \label{eq:row-update}
\end{equation}
Here $i$ is the index of a row other than the pivot row, so $i\ne\rho$,
and $j$ is the column index.
Subtracting the stated multiple of row $\rho$ sets the coefficient in
column $k$ to zero and can create nonzero entries in other columns,
a change called fill-in when a previously zero entry becomes nonzero.
These substitutions also make the coefficient expressions longer by
combining their numerators and denominators, so elimination can require
more time and memory as it proceeds. Redundant equations add to this
work~\cite{Kant}, so Kira uses modular elimination to
select equations before processing their full symbolic coefficients.

\subsection{Removing redundant equations and selecting target dependencies}
To reduce the number of equations it must solve, Kira identifies rows
that follow from other equations. Elimination reduces every coefficient
in such a row to zero, allowing the solver to omit that row. Kira
tests these dependencies with modular arithmetic~\cite{KiraOriginal}
by substituting sample values $\boldsymbol x_0,\epsilon_0$ and mapping
the coefficients to a finite field $\mathbb F_P$ for a prime integer $P$:
\begin{equation}
 A(\boldsymbol x,\epsilon)
 \longmapsto A(\boldsymbol x_0,\epsilon_0)\bmod P.
 \label{eq:modular-selection}
\end{equation}
Integers from $0$ to $P-1$ represent the elements of $\mathbb F_P$, and
addition and multiplication keep results in this range by taking the
remainder upon division by $P$. Division uses the modular inverse of a
nonzero element, which multiplies that element to one modulo $P$, so the
sampling procedure excludes points where a required denominator vanishes.
This arithmetic keeps field elements bounded in size while allowing the
solver to identify independent rows and retain their original symbolic
identities.

The matrix rank is the largest number of independent rows, where no row
is a linear combination of the others in the chosen set. Some sample
values can make otherwise independent rows dependent, giving a rank
smaller than that of the symbolic matrix. Checks must account for this
possibility when using the probabilistic equation-selection method
developed by Kant~\cite{Kant}.

Kira also tracks which independent equations contribute to the target
reductions, using the dependencies of the forward-eliminated system for
this selection in Kira 3~\cite{Kira3}. Seed selection therefore supplies
the integrals from which Kira generates equations, and equation selection
then retains the relations needed for the targets so that the solver can
perform the full reduction on this smaller system.

\subsection{Back substitution and the final symbolic result}
To obtain $C_{\boldsymbol a m}$ from the retained equations, Kira applies
back substitution to the pivot rules produced by forward elimination.
Since these rules can contain other pivot integrals, back substitution
inserts their expressions recursively until the right-hand sides contain
only masters, giving Eq.~\eqref{eq:reduction}. Kira's original symbolic
solver carries out this calculation using rational-function arithmetic
through the computer algebra system Fermat~\cite{KiraOriginal}.

Kira also solves systems at finite-field sample points and reconstructs the
rational coefficients with FireFly, the reconstruction library introduced
into the solver in Kira 2.0~\cite{Kira2}. In this procedure, each sampled
coefficient has the form
\begin{equation}
 C_{\boldsymbol a m}(\boldsymbol x_0,\epsilon_0)\bmod P.
 \label{eq:sampled-coefficient}
\end{equation}
The modular operations use exact arithmetic, and reconstruction recovers
$C_{\boldsymbol a m}(\boldsymbol x,\epsilon)$ from sufficient samples
before checking the result at additional points. Reconstruction therefore
forms part of the calculation of the symbolic reduction coefficients,
after which the master integrals still require evaluation~\cite{Kira2}.

Because seed selection determines the equations that enter this procedure
and elimination decisions determine the substitutions and intermediate
expressions, the proposed ML models would guide these two operations to
reduce the time and memory needed for the reduction.
